\section{Lecture 5 (September 18th)}
\begin{defi}
Notice that the linear fractional transformation defined for $ad-bc\ne 0$ as
\[T(z)=\dfrac{az+b}{cz+d}\]
we see that the derivative is 
\[T'(z)=\dfrac{ad-bc}{(cz+d)^2}\ne 0\]
and that it is analytic and $T'$ is never 0. Meanwhile,
\[T^{-1}(z)=\dfrac{dz-b}{-cz+a}\]
Conformal are by definition, angle preserving maps (both size and orientation).
\end{defi}
\vspace{2ex}
\begin{defi}
Let $z=z(t)$ for $a\leq t\leq b$ represent the contour $C$. Also, let $w=f(z)$, that is, $w(t)=f(z(t))$ represent the contour $\mathrm{\Gamma} $. If $z_0=z(t_0)$ and $w_0=f(z_0)=f(z(t_0))$, $f$ is analytic at $z_0$ such that $f'(z_0)\ne 0$. Then, from \[w'(t_0)=f'(z_0)z'(t_0)\]
we have
\[\mathrm{arg}\,w'(t_0)=\mathrm{arg}\,f'(z_0)+\mathrm{arg}\,z'(t_0)\]
Thus, we see that analytic functions with nonzero deriviatives are conformal mappings.
\end{defi}
\vspace{2ex}
\begin{defi}
Consider
\[T(z)=\dfrac{az+b}{cz+d}\]
(with $ad-bc\ne 0$) and 
\[S(z)=\dfrac{ez+f}{gz+h}\]
(with $eh-gf\ne 0$) we ultimately have
\[(T \circ S)(z)=\dfrac{(ae+bg)z+(af+bh)}{(ce+dg)z+(cf+dh)}\]
also,
\[T^{-1}(z)=\dfrac{dz-b}{-cz+a}\]
Ultimately, we see that these mappings are equivalent to matrix multiplication. 
\end{defi}
\vspace{2ex}
\begin{prop}
(The fixed points of $T(z)=z$) We see that
\[\dfrac{az+b}{cz+d}=z\qquad \implies \qquad cz^2+(d-a)z-b=0\]
and that there can't be more than two fixed points. If $S$ and $T$ are linear fractional transformations such that $S(z_{k})=T(z_{k})$ for $k=1,2,3$, then $S=T$ since $S^{-1}\circ T=I$.
\end{prop}
\vspace{2ex}
\begin{thm}
If distinct points $\{z_1,z_2,z_3\}$ and distinct points $\{w_1,w_2,w_3\}$ are points of $S^2={\bm C}\cup \{\infty \}$ then there exists a unique linear fractional transformation $T$ such that $T(z_{k})=w_{k}$ for $k=1,2,3$.
\[\{a,b,c\}\rightarrow \{0,1,\infty \}\leftrightarrow \{\alpha ,\beta ,\gamma \}\]
Then
\[T(z)=\dfrac{(b-c)(z-a)}{(b-a)(z-c)}\qquad \&\qquad S(z)=\dfrac{(\beta -\gamma )(z-\alpha )}{(\beta -\alpha )(z-\gamma )}\]
\end{thm}
\vspace{2ex}
\begin{ex}
Find $V$ such that
\[\{0,1,2\}\rightarrow \{-1,0,4\}\]
\end{ex}
\vspace{2ex}
\begin{proof}
\[S(z)=\dfrac{-4(z+1)}{1(z-4)}=\dfrac{-4z-4}{z-4}\]
(we say it maps to$\{0,1,\infty \}$) which is equivalent to
\[\begin{pmatrix}
-4&-4\\1&-4
\end{pmatrix}
\]
For $T$ that sends $\{0,1,2\}$ to the same points, we find the matrix, and $V(z)$ is given as
\[V(z)=\dfrac{8z-8}{-3z+8}\]
\end{proof}
\vspace{2ex}
\begin{ex}
Let 
\[T(z)=\dfrac{1+z}{1-z}\]
which corresponds to 
\[\begin{pmatrix}
1&1\\-1&1
\end{pmatrix}
\]
Let $U=\{|z|<1\}$, and say that we want to find $T(U)$. Notice that 
\[\{-1,0,1\}\mapsto \{0,1,\infty \}\]
As $T([-1,1])$ maps to the positive $x$-axis, we can see that $T(-1)=\infty $ and that $T(|z|=1)$ is a line passing $T(-1)=0$. Then, $T(i)=i$ and $T(|z|=1)$ is the $y$-axis. Another way to do this is through conformality of $T(|z|=1)$ and the positive real axis being $\pi /2$ at $z=0$.
\end{ex}
\vspace{2ex}
\begin{defi}
Let $U$ be the unit disk and $T=\{|z|=1\}$ be a unit circle. For $\alpha \in U$, we define
\[\phi _{\alpha }(z)=\dfrac{z-\alpha }{1-\bar{\alpha }z}:\begin{pmatrix}
1&-\alpha \\
-\bar{\alpha }&1
\end{pmatrix}
\]
Then $\phi _{\alpha }$ is analytic on $\bar{U}$ such that $\phi ^{-1}_{\alpha }=\phi _{-\alpha }$, $\phi _{\alpha }(0)=-\alpha $, $\phi_{\alpha }(\alpha )=0$, $\phi _{\alpha }(U)=U$, $\phi _{\alpha }(T)=T$, and $\phi _{\alpha }'(0)=1-|\alpha |^2$, $\phi '(\alpha )=1/(1-|\alpha |^2)$.
\end{defi}
\vspace{2ex}
\begin{proof}
We'll show that $\phi _{\alpha }(T)=T$.
\[\begin{align*}
\phi _{\alpha }(e^{it})=&\dfrac{e^{it}-\alpha }{1-e^{-it}\bar{\alpha }}=\dfrac{e^{it}(1-e^{-it}\alpha )}{1-e^{-it}\bar{\alpha }}\\
|\phi _{\alpha }(e^{it})|=1
\end{align*}\]

\end{proof}
\vspace{2ex}

